\subsection{DIRECTED SETS}

DEFINITION 7. A preordered set E is said to be right directed (resp. left directed) if every subset of two elements of E is bounded above (resp. bounded below).

In place of ``right directed'' we shall often use the expression ``directed with respect to the relation \(\leqslant\) '', and analogous expressions when the preorder relation is denoted by some other sign. For example, if S is a set of subsets of a set A, we say that S is directed with respect to the relation \(\subset\) (resp. \(\supset\) ) if, for each subset \(\{X, Y\}\) consisting of two elements of S, there exists \(Z \in S\) such that \(X \subset Z\) and \(Y \subset Z\) (resp. \(X \supset Z\) and \(Y \supset Z\) ).

Examples

\begin{enumerate}
\def\labelenumi{(\arabic{enumi})}
\tightlist
\item
  An ordered set which has a greatest element is right directed.
\end{enumerate}

* (2) In a topological space, a fundamental system of neighbourhoods of a point is directed with respect to the relation \(\supset\) .

\begin{enumerate}
\def\labelenumi{(\arabic{enumi})}
\setcounter{enumi}{2}
\tightlist
\item
  The set of submodules of finite type of an arbitrary module is directed with respect to the relation \(\subset\cdot_{*}\)
\end{enumerate}

PROPOSITION 10. In a right directed ordered set E, a maximal element \(a\) is the greatest element of E.

For every \(x \in \mathbf{E}\) there exists by hypothesis \(y \in \mathbf{E}\) such that \(x \leqslant y\) and \(a \leqslant y\) ; since \(a\) is maximal, \(y = a\) .

A right directed preordered set is left directed with respect to the opposite ordering. Any product of right directed sets is right directed. On the other hand, a subset of a right directed set is not necessarily right directed. However, a cofinal subset F of a right directed set E is always right directed; for given \(x, y \in \mathbf{F}\) there exists \(z \in \mathbf{E}\) such that \(x \leqslant z\) and \(y \leqslant z\) , and then \(t \in \mathbf{F}\) such that \(z \leqslant t\) .

